\section{Доказательство нижней оценки}\label{app:lower-constants}

\begin{proof}[Доказательство теоремы~\ref{thm:lower}\proofup{thm:lower}]
Достаточно оценить четыре слагаемых $B_n,B_S,B_X,B_Q$
из~\eqref{eq:lower-four-terms}: их сумма имеет порядок $T_0$.
Положим $a=L_1R$ и будем использовать величину $N_0$ из~\eqref{eq:complexity-parameters}.
Числа $G_H$, $\widehat\psi(1)$ и $c_H$ вычислены
в~\eqref{eq:lower-profile-numbers}.

\subsection{Член с \texorpdfstring{$L_0$}{L0}}
По лемме~\ref{lem:lower-pair}
\begin{equation}\label{eq:lower-n-number}
 B_n=n\le9\mathcal T^\ast.
\end{equation}
Если выполнено~\eqref{eq:lower-S-regime}, применима
лемма~\ref{lem:lower-S}. Иначе
\[
 B_S<\sqrt{24/c_H}\,(G_H+6)n<50n.
\]
В обоих случаях
\begin{equation}\label{eq:lower-S-number}
 B_S=\sqrt{\frac{nL_0R^2}{\varepsilon}}\lesssim\mathcal T^\ast.
\end{equation}

\subsection{Члены, получаемые из семейства X}
Используем обозначения $\eta,\rho,D_\delta,U$ из~\eqref{eq:lower-X-parameters}.
По предположению~\ref{asm:accuracy}
$\varepsilon/\delta\le a^2/3$, откуда
\begin{equation}\label{eq:lower-rho-number}
 \eta\le\log(1+2a^2/3)\le\log N_0,\qquad
 1\le\rho\le1+\eta\le1+\log N_0.
\end{equation}
Для второй оценки достаточно неравенств
$z/(1+z)\le\log(1+z)\le z$ при $z=2\varepsilon/\delta$.
Кроме того,
\begin{equation}\label{eq:lower-D-bound}
 D_\delta\sqrt{1+\delta/\varepsilon}\lesssim\log N_0.
\end{equation}
Действительно, $\eta\delta/\varepsilon\le2$ по неравенству
$\log(1+z)\le z$, а
$\sqrt{1+\delta/\varepsilon}\le1+\delta/(2\varepsilon)$. Поэтому
\[
 D_\delta\sqrt{1+\delta/\varepsilon}
 \le\max\left\{\sqrt{\frac{12(\eta+2)}{c_H}},
                 \frac{4(G_H+6)}{c_H}\rho\right\}
 \lesssim\log N_0.
\]

Покажем, что
\begin{equation}\label{eq:lower-Y-number}
 \sqrt n\,U\lesssim\log N_0\,\mathcal T^\ast.
\end{equation}
При $U\ge\sqrt n$ это следует из леммы~\ref{lem:lower-X}
и~\eqref{eq:lower-rho-number}.
При $1\le U<\sqrt n$ та же лемма даёт $U^2\lesssim\rho\mathcal T^\ast$.
Перемножая это неравенство и~\eqref{eq:lower-n-number}, получаем
$\sqrt n\,U\lesssim\sqrt\rho\,\mathcal T^\ast\lesssim\log N_0\,\mathcal T^\ast$.
Наконец, при $U<1$ достаточно $\sqrt n\,U\le n$.
Из $a=7UD_\delta/2$, \eqref{eq:lower-D-bound}
и~\eqref{eq:lower-Y-number} следует
\begin{equation}\label{eq:lower-X-number}
 B_X=\sqrt n\,a\sqrt{1+\delta/\varepsilon}
 \lesssim\log^2N_0\,\mathcal T^\ast.
\end{equation}

\subsection{Член с показателем \texorpdfstring{$2/3$}{2/3}}
Положим $t=(2a/7)^{2/3}$, $m=\lfloor n/2\rfloor$ и
$\eta_m=\log(1+2\varepsilon m/\delta)$.
По предположению~\ref{asm:accuracy}
\[
 \eta_m\le\log(1+na^2/3)\le2\log N_0.
\]
Если выполнены условия~\eqref{eq:lower-Q-regime},
лемма~\ref{lem:lower-Q} даёт
\begin{equation}\label{eq:lower-Q-number}
 B_Q=na^{2/3}\lesssim(1+\eta_m)^2\mathcal T^\ast
 \lesssim\log^2N_0\,\mathcal T^\ast.
\end{equation}
Если хотя бы одно условие нарушено, оценка получается из уже доказанных:
\begin{itemize}
\item При $a<21/2$ имеем $B_Q<5n$.
\item Если $a\ge21/2$, но $m<4t$, то $t>1$ и
$n\le2m+1<9t$. Поэтому
\[
 \frac{B_Q}{\sqrt n\,a}=\frac{\sqrt n}{a^{1/3}}
 <3(2/7)^{1/3}<2,
 \qquad B_Q<2B_X.
\]
\item Если $24\eta_m>\widehat\psi(1)t$, то
\[
 B_Q<(7/2)^{2/3}\frac{24}{\widehat\psi(1)}n\eta_m
 \lesssim n\log N_0.
\]
\end{itemize}
Таким образом, оценка~\eqref{eq:lower-Q-number} верна во всех случаях.
Складывая оценки для четырёх слагаемых, получаем
\[
 T_0\sim B_n+B_S+B_X+B_Q
 \lesssim\log^2N_0\,\mathcal T^\ast.
\]
\end{proof}
